% ============================================================================
% Section 1: Memory-Mapped GPIO (Task 1)
% Handout requires: peripheral addresses and offsets, bit masks used,
% short explanation of volatile, Task 1 oscilloscope capture.
% ============================================================================
\section{Memory-Mapped GPIO}
\label{sec:gpio}

PC13 was configured as a push-pull output and toggled using explicit
\reg{volatile} pointers built from addresses in RM0091~\cite{rm0091}, without
any HAL GPIO function. Each register address is the peripheral base address plus
the register offset, $A = B + O$ (\cref{tab:gpio-addr}). The build checks every
address against the CMSIS device header with \reg{\_Static\_assert}.

% Values below are taken from Core/Src/task1_gpio.c. The red "Source" entries
% are suggested locations - check each one in your copy of RM0091 and correct
% the section/table number if your revision differs.
\begin{table}[htbp]
  \centering
  \caption{Base addresses, offsets and resulting register addresses used in Task 1.}
  \label{tab:gpio-addr}
  \begin{tabular}{@{}lllll@{}}
    \toprule
    \textbf{Register} & \textbf{Base $B$} & \textbf{Offset $O$} & \textbf{Address $B+O$} & \textbf{Source} \\
    \midrule
    \reg{RCC\_AHBENR}  & \texttt{0x40021000} & \texttt{0x14} & \texttt{0x40021014} & \fillin{RM0091 \S2.2.2 Tab.~1; \S6.4.6} \\
    \reg{GPIOC\_MODER} & \texttt{0x48000800} & \texttt{0x00} & \texttt{0x48000800} & \fillin{RM0091 \S8.4.1} \\
    \reg{GPIOC\_ODR}   & \texttt{0x48000800} & \texttt{0x14} & \texttt{0x48000814} & \fillin{RM0091 \S8.4.6} \\
    \reg{GPIOC\_BSRR}  & \texttt{0x48000800} & \texttt{0x18} & \texttt{0x48000818} & \fillin{RM0091 \S8.4.7} \\
    \reg{GPIOC\_BRR}   & \texttt{0x48000800} & \texttt{0x28} & \texttt{0x48000828} & \fillin{RM0091 \S8.4.11} \\
    \bottomrule
  \end{tabular}
\end{table}

% Masks below match the current implementation in task1_gpio.c. Update this
% table if you change the code (e.g. switch the toggle to BSRR/BRR).
\begin{table}[htbp]
  \centering
  \caption{Bit masks used in Task 1.}
  \label{tab:gpio-masks}
  \begin{tabularx}{\linewidth}{@{}llllX@{}}
    \toprule
    \textbf{Register} & \textbf{Field} & \textbf{Mask} & \textbf{Operation} & \textbf{Purpose} \\
    \midrule
    \reg{RCC\_AHBENR}  & IOPCEN (bit 19)          & \texttt{0x00080000} & \texttt{|=}  & Enable the GPIOC clock; other clock enables kept \\
    \reg{GPIOC\_MODER} & MODER13[1:0] (bits 27:26) & \texttt{0x0C000000} & \texttt{\&= \textasciitilde} & Clear PC13's mode field; other pins unchanged \\
    \reg{GPIOC\_MODER} & MODER13 = \texttt{01}     & \texttt{0x04000000} & \texttt{|=}  & General-purpose output mode \\
    \reg{GPIOC\_ODR}   & ODR13 (bit 13)           & \texttt{0x00002000} & \texttt{|=}  & Known starting level (high) \\
    \reg{GPIOC\_ODR}   & ODR13 (bit 13)           & \texttt{0x00002000} & \texttt{\^{}=} & Toggle PC13 every half period \\
    \bottomrule
  \end{tabularx}
\end{table}

\paragraph{Choice of output register.}
\writeup{Why you toggled with ODR (read-modify-write) rather than BSRR/BRR (single atomic write), and why that is safe in this program. 2--3 sentences.}

\paragraph{Why \texttt{volatile}.}
\writeup{Short explanation: the register can change, or react to an access, independently of program flow, so the compiler must perform every read and write rather than caching or removing them. Give one concrete example of what goes wrong without it (e.g. a repeated write or a polling loop being optimised away).}

\paragraph{GPIO clock disabled.}
\writeup{One or two sentences: what the board does if IOPCEN is not set (GPIOC is unclocked, register writes have no effect, PC13 stays in its reset state and never toggles).}

\scopefig{task1_pc13.png}{PC13 square wave generated by memory-mapped GPIO
  (\reg{task1\_half\_period\_ms} $=$ \fillin{5}~ms). Measured period
  \fillin{\ldots~ms}, frequency \fillin{\ldots~Hz}.}{fig:task1}
